% paper4_tillage.tex -- Paper 4. Typeset alone with this file, or inside
% the volume through main.tex.
\documentclass[../main.tex]{subfiles}
\begin{document}
\begin{paper}{Tillage and Residue Management Effects on Soil Organic Carbon and Yield}{Tillage, residues and soil organic carbon}
\label{paper:tillage}

\begin{paperabstract}
Reduced tillage and the retention of crop residues are promoted to protect
soils and build soil organic carbon, but their effect on crop yield varies
with climate and management. This paper reports a factorial trial combining
two tillage systems, conventional and no-till, with two residue treatments,
removed and retained, at three sites. Soil organic carbon in the topsoil is
compared with its value at the start of the trial, and grain yield is compared
across the four combinations. The data are synthetic; the paper shows how a
two-factor trial is analysed and how changes in soil organic carbon are
reported alongside yield.
\end{paperabstract}
\paperkeywords{conservation agriculture; no-till; crop residues; soil organic carbon; factorial trial}

\section{Introduction}

Soil organic carbon lies at the centre of soil health\index{soil health}. It supports soil
structure, water retention and nutrient supply, and agricultural soils that
have lost carbon through cultivation have a large capacity to regain it.
Practices that increase the soil carbon pool include reduced tillage, the
retention of residues, cover crops and the addition of manure\index{organic amendments!manure}s, and the
resulting gains in soil carbon can support higher crop yields on degraded
land \parencite{lal2004}.

Conservation agriculture\index{conservation agriculture} combines three principles: minimal soil disturbance,
permanent soil cover and crop diversification. A global meta-analysis of
paired yield observations found that no-till\index{tillage!no-till} on its own tends to reduce
yields relative to conventional tillage, but that the reduction is smallest
when no-till is combined with residue retention\index{crop residues} and crop rotation, and that
the combination can raise rainfed yields in dry climates
\parencite{pittelkow2015}. The yield effect of reduced tillage therefore
depends on whether the other principles are applied with it.

Whether these principles can be applied in practice is a separate question.
In smallholder farming systems, crop residues are often needed as livestock
feed, weeding without tillage can increase labour, and the evidence for
benefits has been argued to be less consistent than its promotion suggests
\parencite{giller2009}. Trials that separate the effect of tillage from the
effect of residues help to show which part of the package produces which
effect.

This paper reports such a trial. Soil organic carbon is determined by wet
oxidation \parencite{walkley1934}. The aims are to estimate the separate and
combined effects of tillage and residue management on topsoil organic carbon\index{soil organic carbon}
and on grain yield, and to test whether the effect of tillage depends on the
residue treatment.

\input{\subfix{methods/paper4_methods}}

\input{\subfix{results/paper4_results}}

\input{\subfix{discussion/paper4_discussion}}

\end{paper}
\end{document}
